\documentclass[11pt,a4paper]{article}
\usepackage[T1]{fontenc}
\usepackage[utf8]{inputenc}
\usepackage{lmodern,microtype,amsmath,amssymb,amsthm,booktabs,array,graphicx}
\usepackage[margin=25mm]{geometry}
\usepackage[colorlinks=true,linkcolor=blue,citecolor=blue,urlcolor=blue]{hyperref}
\usepackage[numbers]{natbib}
\newtheorem{theorem}{Proposition}
\newtheorem{lemma}[theorem]{Lemme}
\newtheorem{assumption}[theorem]{Hypothèse}
\newcommand{\code}[1]{\nolinkurl{#1}}
\newcommand{\Com}{\operatorname{Com}}
\newcommand{\Zq}{\mathbb Z_q}
\hypersetup{pdftitle={Horizon : validation correlee de la relation de plages},pdfauthor={Dossier technique ZK-agentiK}}
\title{Horizon : validation corrélée de la relation de plages\\
  \large Supplément algébrique et protocole expérimental}
\author{Dossier technique du dépôt ZK-agentiK}
\date{17 octobre 2026 \quad Version expérimentale 1}
\begin{document}
\maketitle
\begin{abstract}
Ce supplément relie la comparaison entière d'un montant et d'un plafond aux
deux engagements dérivés vérifiés par Horizon. Il donne une démonstration
conditionnelle et un protocole de corrélation entre calcul JavaScript BigInt,
oracle MATLAB/GNU Octave et exécution réelle de Bulletproofs en Rust. Les
conditions de succès distinguent refus applicatif et rejet cryptographique.
Des essais finis, même concordants, ne prouvent ni l'absence de défaut, ni la
sécurité de la primitive. La source numérique est portable ; le runtime
effectivement exécuté doit être lu dans chaque artefact. Aucune exécution
MATLAB n'est déduite d'un résultat Octave. Le profil n'est pas post-quantique.
\end{abstract}

\section{Énoncé : relation arithmétique et politique}
Le profil examiné est \code{horizon-range-v1}, dans
\code{labs/horizon/proof/src/lib.rs}. Ce supplément accompagne le manuscrit
\code{docs/whitepaper/whitepaper.tex} sans le remplacer. La base de la campagne
est le commit \code{5cf0ee2} ; chaque exécution doit identifier la révision
effective de ses sources et ses éventuelles modifications locales.

Posons $U=2^{64}-1$ et $N=2^{32}$. Le domaine d'entrée est
$\mathcal D=\{0,\ldots,U\}^2$. La décision numérique est
\begin{equation}\label{eq:predicate}
 P(a,c)\ \Longleftrightarrow\ 1\le a\le N\quad\land\quad
 a\le c\quad\land\quad c-a<N.
\end{equation}
Cette relation n'est pas simplement $a\le c$. Par exemple $(0,0)$,
$(N+1,N+1)$ et $(1,N+1)$ satisfont la comparaison mais pas $P$.
Le plafond $c$ est un \code{u64} et peut atteindre $2^{33}-1$ ; imposer
$c\le2^{32}-1$ changerait la relation.

\paragraph{Frontière de politique.}
Une acceptation cryptographique exige aussi les signatures, la jointure du
dossier, la devise EUR, l'exposant 2, les clés et rôles épinglés, les dates,
le challenge indépendant et le mode permis. L'autorisation d'un effet exige
encore révocation, époque, nonce et opération métier dans le contrôleur.
L'oracle numérique ne reçoit aucun de ces éléments. Il ne peut donc pas
prononcer une autorisation métier. Le plafond reste par dossier : deux montants
de 60 sous un plafond de 100 n'établissent pas un budget cumulé de 100 ;
une franchise, une conversion d'unité et une règle de paiement seraient
des relations supplémentaires.

\section{Du witness entier aux engagements}
Soit $\mathbb G$ le groupe Ristretto d'ordre premier
\[
 q=2^{252}+27742317777372353535851937790883648493.
\]
Les encodages canoniques et l'abstraction de groupe sont décrits dans
RFC~9496 \cite{ristretto}. Les bases $G,H$ sont les paramètres publics fixes
de \code{PedersenGens::default()}, jamais choisis par l'appelant. Deux
autorités construisent les engagements signés
\[
 A=aG+r_aH,\qquad C=cG+r_cH,\qquad r_a,r_c\in\Zq.
\]
Le vérifieur dérive les points, sans accepter ceux proposés séparément par
le prover :
\begin{equation}\label{eq:derived}
 X=A-G,\qquad D=C-A.
\end{equation}
Le witness de la preuve de plages agrégée est
\[
 x=a-1,\quad d=c-a,\quad r_X=r_a,\quad r_D=r_c-r_a\pmod q.
\]
Horizon utilise \code{RangeProof} de Bulletproofs 5.0.0 avec deux valeurs et
32 bits \cite{rangeapi}. Ce chemin est une preuve spécialisée de plages ;
il ne compile pas la politique entière en un circuit R1CS universel.
Le mot « circuit » dans la campagne désigne cette relation cryptographique.

\begin{lemma}[Équivalence entière]
Pour $(a,c)\in\mathcal D$, $P(a,c)$ équivaut à l'existence d'entiers
$x,d\in[0,N-1]$ tels que $a=x+1$ et $c=x+1+d$.
\end{lemma}
\begin{proof}
Si $P$ est vraie, $x=a-1$ et $d=c-a$ satisfont les bornes. Réciproquement,
$a=x+1\in[1,N]$ et $c-a=d\in[0,N-1]$. Ces identités donnent aussi
$c\le2N-1$, sans addition susceptible de dépasser $U$.
\end{proof}

\begin{lemma}[Complétude arithmétique]
Si $P(a,c)$ est vraie, les ouvertures précédentes produisent exactement
les points de \eqref{eq:derived}.
\end{lemma}
\begin{proof}
$xG+r_XH=(a-1)G+r_aH=A-G$. De même,
$dG+r_DH=(c-a)G+(r_c-r_a)H=C-A$. Les deux valeurs sont dans les plages
requises ; la complétude cryptographique demeure celle de la primitive.
\end{proof}

\section{Correction conditionnelle : connaissance et liaison}
\begin{assumption}[Émission, primitive et implémentation]
Les autorités ont construit ou contrôlé les ouvertures de $A$ et $C$ depuis
des valeurs du domaine \code{u64}, avec les métadonnées appropriées.
Le logarithme discret entre $G$ et $H$ n'est pas connu d'un adversaire
efficace. La primitive fournit l'argument de connaissance des ouvertures
dans les plages sous son modèle de sécurité. Encodages, paramètres,
transcript, RNG et vérification implémentent fidèlement cette primitive.
\end{assumption}
Le papier Bulletproofs et sa révision corrigée sont la référence, sans
réduction ni audit logiciel réinventé ici \cite{bulletproofs}.
Une simple existence mathématique d'ouvertures serait insuffisante : puisque
$H$ engendre $\mathbb G$, tout point admet une ouverture à toute valeur
scalaire pour un masque approprié. Le passage à une valeur attestée exige
la connaissance extractible et la liaison calculatoire.

\begin{lemma}[Deux ouvertures incompatibles]
Si $vG+rH=v'G+r'H$ et $v\not\equiv v'\pmod q$, alors
\[
 \log_G H=(v-v')(r'-r)^{-1}\pmod q.
\]
\end{lemma}
\begin{proof}
L'égalité implique $(v-v')G=(r'-r)H$. Si $r'=r$, les valeurs seraient
congrues ; sinon $r'-r$ est inversible dans le corps d'ordre premier.
La formule donne la relation discrète supposée inconnue.
\end{proof}

\begin{theorem}[Liaison aux entiers attestés, sans wrap]
Supposons les hypothèses précédentes. Une preuve acceptée dont on extrait
$(x,d,r_X,r_D)$ avec $0\le x,d<N$ impose $P(a,c)$ pour les valeurs
attestées, sauf échec d'extraction ou rupture de liaison.
\end{theorem}
\begin{proof}
Les points dérivés donnent une ouverture de $A$ à $a'=x+1$ avec masque
$r_X$, et de $C$ à $c'=x+1+d$ avec masque $r_X+r_D$. Si ces valeurs
diffèrent modulo $q$ des ouvertures attestées, le lemme précédent donne
une rupture de liaison. Sinon $a\equiv a'$ et $c\equiv c'\pmod q$.
Or $0\le a,c\le U<q$ et $0\le a',c'\le2^{33}-1<q$.
Deux représentants dans $[0,q)$ congrus modulo $q$ sont égaux.
Donc $a=a'$ et $c=c'$, et l'équivalence entière conclut.
\end{proof}
La borne $U<q$ est essentielle : une autorité attestant un entier sans
domaine pourrait signer une valeur $a'+q$ ayant le même scalaire. De même,
la signature d'un point fourni aveuglément ne garantit pas que le point
représente le montant issu du système métier. Ces risques ne sont pas
supprimés par une concordance de tests.

Les signatures d'autorité sont vérifiées hors preuve de plages. Le mode A
inclut le challenge complet dans le transcript ; B réutilise une preuve
historique et signe une nouvelle présentation sous la clé du détenteur.
B ne démontre pas une connaissance fraîche du witness. La composition
suppose aussi l'inforgeabilité des signatures, la résistance aux collisions
du digest, la séparation des domaines et la comparaison complète du challenge.
Aucune probabilité numérique de sécurité n'est déduite des essais.

\section{Oracle MATLAB/Octave : exactitude des entiers}
L'entrée de \code{math/circuit/circuit_oracle.m} est un tableau JSON de
1 à 10\,000 objets contenant exactement \code{id}, \code{category},
\code{amount} et \code{cap}. Les valeurs numériques sont des chaînes
décimales canoniques, jamais des nombres JSON convertis implicitement en
\code{double}. L'API \code{jsondecode} documente cette conversion et la
normalisation des noms de champs \cite{matlabjson}. Des contrôles de forme,
de clés, de doublons et d'encodage complètent donc le décodage standard.
Un champ \code{expected} est refusé. L'oracle ne lit aucune réponse attendue
dans le corpus.

\begin{lemma}[Conversion décimale sans saturation]
Posons $Q=\lfloor U/10\rfloor=1844674407370955161$ et $R=5$.
Pour un préfixe $v\in[0,U]$ et un chiffre $b\in[0,9]$, l'opération
$10v+b$ est représentable si et seulement si $v<Q$, ou $v=Q$ et $b\le R$.
\end{lemma}
\begin{proof}
$U=10Q+R$. Si $v<Q$, $10v+b\le10Q-1\le U$. Si $v=Q$, la condition
est exactement $b\le R$. Si $v>Q$, $10v+b\ge10Q+10>U$.
L'induction sur les chiffres justifie le contrôle avant multiplication et
addition. $Q$ est construit depuis de petits blocs exactement représentables,
sans conversion d'un grand littéral flottant.
\end{proof}

Le calcul de $x$ est précédé de $a\ge1$ et celui de $d$ de $c\ge a$.
Cette garde est indispensable : les opérations entières peuvent saturer
aux limites des types dans MATLAB et Octave \cite{matlabintegers,octaveintegers}.
Une différence négative est rendue par la chaîne vide, pas par zéro ; les
autres différences sont rendues exactement, y compris pour un cas refusé.

Le formatage décimal évite aussi la division entière au bord du type.
Pour une valeur supérieure à $2^{32}-1$, une division binaire par 10 garde
un reste dans $[0,9]$ ; avant réduction, ce reste est au plus 19. Le
quotient est construit bit par bit. Pour une valeur au plus $2^{32}-1$, le
chemin court utilise \code{floor(double(v)/10)} : l'entier est représenté
exactement et l'erreur d'arrondi du quotient est très inférieure à $0{,}1$,
distance minimale d'un quotient non entier au prochain entier. Aucun
grand \code{uint64} n'est envoyé dans ce chemin.

\paragraph{Observation de développement.}
Sur GNU Octave 7.3.0 dans l'image locale enregistrée, une première version
utilisant \code{idivide} au maximum du type n'a pas satisfait les tests de
formatage : $2^{64}-1$ a été rendu avec une terminaison décimale erronée.
Ces essais ont échoué avant toute publication d'un résultat positif.
La source livrée n'utilise pas cette division. Ce constat est un résultat
de runtime local, pas une conclusion sur toutes les versions MATLAB/Octave.

\section{Protocole de corrélation des trois chemins}
Pour chaque vecteur $i$, notons $J_i$ la décision BigInt, $M_i$ celle du
runtime numérique réellement exécuté et $R_i$ le résultat Rust approprié.
Le corpus et chaque résultat doivent porter des identifiants uniques et
être reliés par leurs empreintes. L'égalité des ensembles d'identifiants,
des valeurs et des effectifs est vérifiée \emph{avant} celle des décisions.
Une omission, duplication, erreur de processus ou sortie illisible ne vaut
jamais un refus métier réussi.

\begin{center}
\fbox{\begin{minipage}{.89\linewidth}\centering
Corpus synthétique commun : identifiants, catégories, entiers décimaux\\[5pt]
$\swarrow\qquad\qquad\downarrow\qquad\qquad\searrow$\\[5pt]
BigInt\quad\textbar\quad MATLAB ou Octave\quad\textbar\quad Rust + Bulletproofs\\[5pt]
$\searrow\qquad\qquad\downarrow\qquad\qquad\swarrow$\\[5pt]
Jointure exacte, décisions comparées, étapes et erreurs conservées
\end{minipage}}
\end{center}
MATLAB et Octave exécutant la même source ne sont pas deux implémentations
indépendantes du prédicat. Des chemins concordants peuvent partager une
erreur de spécification. La preuve algébrique et les contre-exemples métier
restent donc nécessaires.

\begin{center}\small
\begin{tabular}{p{.24\linewidth}p{.66\linewidth}}
\toprule
Classe d'essai & Condition de réussite et sens du résultat \\
\midrule
Cas arithmétique valide & $J_i=M_i=\mathrm{vrai}$ ; preuves A et B générées,
 puis acceptées par un processus vérifieur distinct sur les attentes épinglées. \\
Cas arithmétique invalide & $J_i=M_i=\mathrm{faux}$ ; le wrapper Rust peut
 refuser avant génération. Ce refus n'est pas un test du vérifieur de plages. \\
Contournement du garde-fou du prover & Construire avec l'API de la bibliothèque
 une tentative hors plage ; observer séparément erreur de génération ou rejet
 cryptographique. Ne pas compter une absence de preuve comme une preuve rejetée. \\
Preuve modifiée & Garder un encodage structurel valide et, lorsque requis,
 re-signer le bundle de fixture. Le refus doit atteindre l'étape cryptographique
 revendiquée, au lieu de s'arrêter à une signature extérieure incorrecte. \\
Contexte ou politique modifiés & Refus d'autorisation distinct d'un échec
 arithmétique ; compte d'effets inchangé dans les essais applicatifs concernés. \\
\bottomrule
\end{tabular}
\end{center}

Le vérifieur de campagne ne reçoit ni masques ni clé privée du détenteur.
Les montants du corpus sont synthétiques et publics pour comparer les
oracles ; cela ne constitue pas un test d'anonymat. Les preuves, métadonnées
et engagements restent corrélables. Un processus séparé démontre une
frontière d'entrée testée, pas une isolation matérielle ou un audit externe.

\section{Couverture, artefacts et interprétation}
Les tailles 128, 1\,024 et 4\,096 sont des paramètres de campagne, pas des
résultats présumés. Un rapport doit donner les cas effectivement terminés,
leurs catégories, les deux verdicts arithmétiques, le point exact de refus
Rust et les preuves réellement vérifiées. Les frontières comprennent
$0$, $1$, $N-1$, $N$, $N+1$, $2N-1$, $2N$, $2^{53}\pm1$, $U$ et les
relations $c<a$, $c=a$, $c-a=N-1$, $c-a=N$. Le parsing teste notamment
$U+1$, les zéros initiaux, signes, nombres JSON, champs supplémentaires et
identifiants dupliqués.

L'autotest local exécuté sous Octave 7.3.0 le 17 octobre 2026 a vérifié
11 cas de frontière, rejeté 36 entrées invalides et refusé un écrasement
de fichier. Il a parcouru les 4\,096 paires du carré $[0,63]^2$, dont
2\,016 satisfont $P$. Cet autotest est numérique ; il n'exécute pas Rust.
Le statut d'une campagne corrélée complète est déterminé par ses artefacts,
et non par ces nombres ou par la compilation du présent document.

Les fichiers publics prévus comprennent \code{vectors.json},
\code{octave.json} ou \code{matlab.json}, \code{prover-results.json},
\code{verifier-results.json}, \code{correlation.json}, les requêtes publiques
de vérification compressées, un rapport HTML et ses figures. Ils sont placés
dans un nouveau répertoire d'exécution ; aucune archive antérieure n'est
écrasée. Les champs \code{runtime}, \code{runtime_version} et
\code{executed_in_matlab} distinguent les moteurs réellement utilisés.
Une exécution GitHub Actions ne remplace pas cette attribution ; une action
configurée ou un job échoué ne vaut pas une exécution MATLAB réussie.

Les figures doivent dériver de résultats présents et porter effectifs,
unités et erreurs. Un panneau sans mesure est marqué « non mesuré ».
Les latences locales n'établissent ni SLA ni supériorité à Vega : les
relations, contrats d'émission et conditions diffèrent. Aucun résultat GPU
n'est déduit de ce test CPU. Des tailles de corpus croissantes ne mesurent
pas la probabilité cryptographique de soundness et ne démontrent pas une
couverture exhaustive du domaine $\mathcal D$.

Les preuves conditionnelles ici sont rédigées et relues, sans assistant
de preuve formelle. Les hypothèses de connaissance, de liaison et de
signatures restent héritées ; l'autorisation distante exactement une fois,
la vérité des sources, les canaux auxiliaires et la qualification bancaire
restent hors portée. Un manifeste SHA-256 lie les octets sous une référence
de confiance ; sa date locale n'est pas un horodatage tiers et son contenu
ne prouve pas la sémantique d'une assertion.

\bibliographystyle{plainnat}
\bibliography{circuit-validation}
\end{document}
